\section{Introduction}\label{sec:introduction}

Volume does not capture all restrictions on how one body can be embedded
in another by a symplectic map. Symplectic capacities measure such
restrictions, giving a way to compare symplectic size with ordinary volume.
This thesis investigates that comparison for four-dimensional convex
bodies. We develop and use algorithms, numerical experiments and proofs
to study which shapes have large capacity and which deformations can
increase it. AI agents participated in this research, and the project also
examines what it means to delegate mathematical work to them.

We use the Ekeland--Hofer--Zehnder capacity $c_{\rm EHZ}$. For a smooth
convex body it is the least action of a closed characteristic on the
boundary; the preliminaries define these trajectories and their extension
to polytopes. Capacity scales like area, so in four dimensions its square
can be compared with volume through the \emph{systolic ratio}
\[
 \operatorname{sys}(K):=
 \frac{c_{\rm EHZ}(K)^2}{2\operatorname{vol}_4(K)}.
\]
A ball has ratio one. Viterbo's conjecture~\cite{Viterbo2000} asserted that
no convex body has a larger ratio. The bound holds for the Gromov width,
which measures the largest symplectic ball that fits inside a body, simply
by volume preservation. Its failure for another capacity therefore shows
that these notions of symplectic size differ even on convex bodies.

Haim--Kislev and Ostrover found precisely such a failure: a product of two
regular pentagons, placed in the complementary $q$- and $p$-planes, has
systolic ratio $(3+\sqrt5)/5>1$~\cite{HaimKislevOstrover2024}.
We call this the HKO body, $K_{\rm HKO}$; its precise position is given in
Section~\ref{sec:hko-local-maximum}. The counterexample does not identify
the largest possible ratio or explain which shapes make it large. Can the
HKO body be improved by a small deformation? What geometric properties
help identify other bodies with large ratios? These are the main questions
pursued here.

The comparison has consequences for both embeddings and convex geometry.
In four dimensions, EHZ capacity coincides with cylindrical capacity,
which measures the smallest symplectic cylinder into which a body embeds
\cite[Corollary~1]{AbbondandoloEdtmairKang2024}. On products of centrally
symmetric bodies with their polars, the volume--capacity bound is equivalent
to the symmetric Mahler volume-product bound~\cite{AKO2014}.
Section~\ref{subsec:preliminaries-capacity-context} explains these links.

\subsection{Computational and empirical approach}

The definition by closed trajectories does not usually give a capacity
value directly. For polytopes, Haim--Kislev's formula~\cite{HK2017}
turns the problem into a finite quadratic program. We give an expanded
derivation of this formula and an implementation for evaluating capacities.
For products of planar polygons, a reduction to at most three facet
directions per factor makes the calculation smaller. A second algorithm
uses the passage geometry of Chaidez and Hutchings~\cite{CH2021}; we prove
its correctness under explicit regularity hypotheses in exact real
arithmetic.

These finite methods make it possible to investigate more than a few
explicitly solvable examples. We sample polytopes, compare geometric
measurements with computed systolic ratios, and test searches that use
capacity evaluations to choose further shapes. Such experiments can reveal
patterns worth explaining and test whether a proposed geometric guide is
useful in a search. The finite formula also supports proofs: feasible
choices in the capacity problem give upper bounds whose variation can
control nearby bodies. Numerical exploration and local analysis thus
address different parts of the capacity--volume question.

\subsection{Results and their scope}

The main local result proves that the HKO body cannot be improved by a
sufficiently small deformation among polytopes with exactly ten facets
(Theorem~\ref{thm:hko-ten-facet-local-maximum}). Equality occurs only under
translations, positive dilations and linear symplectic maps, all of which
preserve the ratio. This gives local optimality at a body with corners
whose ratio exceeds the ball's. For comparison, Abbondandolo and Benedetti
prove local systolic maximality for Zoll contact forms, including the
volume--capacity inequality near the ball~\cite{AB2023}. Our theorem
fixes the number of facets and gives no conclusion for perturbations that
add facets, such as the hyperplane cuts studied by Haim--Kislev
\cite{HaimKislev2025Zoll}. Section~\ref{subsec:hko-related-local-results}
gives the detailed comparison.

There are other local maxima in the same ten-facet space. A rational
hexagon--quadrilateral product has ratio one and a non-strict local
maximum, with a two-parameter family of equality cases after removing the
same symmetries (Theorem~\ref{thm:non-hko-ten-facet-local-maximum}). This
is not another counterexample to Viterbo's inequality; it shows that local
maximality in this space is not confined to the HKO body.

Results for particular product families give further restrictions on how
a large ratio can arise. Independently stretching or shearing the two
regular-pentagon factors cannot improve the HKO ratio, even for large
deformations (Theorem~\ref{thm:pentagon-affine-capacity}); the rotation
family has an explicit capacity formula
(Theorem~\ref{thm:rotated-pentagon-product-formula}). If one factor is a
regular pentagon and the other is any centrally symmetric planar convex
body, the sharp upper bound is $(10+6\sqrt5)/25<1$
(Theorem~\ref{thm:pentagon-symmetric-partner}), so symmetry of the partner
rules out a counterexample in this class. The placement of the factor
planes matters too: every product in symplectically orthogonal planes
admits a small rigid rotation that increases capacity at fixed volume
(Theorem~\ref{thm:product-initial-rotation}).

The empirical study finds a strong inverse rank association between the
systolic ratio and the total symplectic area of the two-dimensional faces,
normalized by $\sqrt{\operatorname{vol}_4}$. Bodies with smaller values
of this measurement tend to have larger ratios in the sampled populations.
The association weakens among low-area or high-ratio bodies. It also
weakens after balancing the shapes of polygon factors, while exact
rotation formulas for pentagon and triangle products give opposite
orderings of area and ratio. These results constrain an explanation of
the observed pattern: it is not a universal rule for increasing capacity.
In the tested local-search comparison, retaining previously active
branches of the capacity calculation improved the achieved ratios, but
endpoint probes still found missed directions of improvement. None of the
random samples or search comparisons reported here produced a numerical
ratio above one (Section~\ref{sec:black-box-datascience}).

\subsection{AI in the research process}

AI use extended to choosing methods and developing mathematical arguments,
software and experiments. In the local maximality work, an agent developed
the argument and exact-computation code from an incomplete suggestion about
controlling nearby bodies by several upper bounds. This provides a
concrete example of delegating mathematical work without prescribing its
solution. Producing readable exposition and keeping agents directed toward
a completed task proved harder. Section~\ref{sec:use-of-ai} reflects on
this experience and on the human judgment needed to direct and assess the
work. Its observations and hypotheses concern this project; the separate
AI-use disclosure records provenance and the extent of AI use.

\subsection{Outline}

Sections~\ref{sec:preliminaries}--\ref{sec:quadratic-program-algorithm-hk2019}
develop the capacity formula, Section~\ref{sec:flow-graph-algorithm-ch2021}
the flow-graph algorithm, and Section~\ref{sec:first-order-perturbations}
the first variation. Section~\ref{sec:hko-local-maximum} proves the local
maximality results without using the flow-graph or data chapters.
Section~\ref{sec:black-box-datascience} presents the empirical study, and
Sections~\ref{sec:rotated-regular-polygons}--\ref{sec:product-rotation}
treat the product families. The remaining sections discuss visualization,
numerics, code and data, AI use, and conclusions. The appendices give
the empirical methods and the exact HKO computation.
